\chapter{Conclusion and Future Research Directions}
\label{chap:conclusion}

\section{Summary of Research Findings}
As Internet of Things (IoT) ecosystems become the backbone of critical infrastructure, industrial automation, and smart urban environments, the vulnerability of static machine learning security models to emerging zero-day threats poses an existential risk. This thesis investigated the fundamental limitations of closed-world supervised Network Intrusion Detection Systems (NIDS) and formulated an autonomous, resource-efficient Continual Learning framework tailored for edge IoT environments.

Through systematic empirical benchmarking on the real-world TON\_IoT intrusion dataset, this research yielded several seminal findings:
\begin{enumerate}
    \item \textbf{Validation of the Zero-Day Blindness:} Conventional supervised tree ensembles fail completely ($F1 = 0.000$) when deployed against unobserved attack vectors, validating the theoretical detection deficit and underscoring the urgent necessity of lifelong learning in network security.
    \item \textbf{Debunking the Full Retraining Ceiling:} Full joint retraining—traditionally considered the optimal upper bound in continual learning—was shown to suffer from severe class imbalance when assimilating compact zero-day attacks, achieving an F1-score of only $0.6483$ on novel MITM traffic. By contrast, our resource-efficient replay framework achieved an F1-score of $0.9773$, demonstrating that balanced historical rehearsal is fundamentally superior to brute-force retraining.
    \item \textbf{Uncovering the Volumetric Flaw of Static Buffers:} While a fixed static replay buffer ($k=100$ samples/class) provides near-perfect storage reduction ($>99.2\%$), it is acutely vulnerable to high-volume attacks (e.g., XSS and Backdoor). When novel attack volume substantially exceeds buffer capacity, maximum catastrophic forgetting spikes to nearly $40\%$.
    \item \textbf{Empirical Triumph of Dynamic Allocation:} Our proposed volume-aware allocation rule, $k_{dynamic} = \max(k_{min}, \lfloor M / K \rfloor)$, dynamically neutralizes class imbalance during tree splitting, slashing maximum catastrophic forgetting by $35.8\%$ while sustaining extreme memory reductions ($96.8\%$ to $99.3\%$) and sub-0.085 second execution latencies.
\end{enumerate}

\section{Research Questions Revisited}
The thesis systematically answered the five foundational research questions posed in Chapter~\ref{chap:introduction}:
\begin{itemize}
    \item \textbf{Answer to RQ1 (Supervised Failure Mode):} Supervised Random Forests achieve an exact $0.000$ detection F1-score on held-out attack classes, confirming that closed-world models are fundamentally incapable of recognizing zero-day threats without continual adaptation.
    \item \textbf{Answer to RQ2 (Anomaly Detection Utility):} Isolation Forests effectively flag out-of-distribution traffic with linear computational complexity, but attain modest recall ($18.86\%$ on MITM) and high false alarm rates, confirming their utility strictly as front-end screening filters rather than standalone classifiers.
    \item \textbf{Answer to RQ3 (Replay Memory Sensitivity):} Empirical sensitivity analysis established that a minimal buffer floor of $k=100$ samples per class ($900$ samples total) provides the optimal Pareto balance, reducing average forgetting by over $58\%$ relative to ultra-small buffers ($k=25$).
    \item \textbf{Answer to RQ4 (Volumetric Attack Dynamics):} Static buffers fail under volumetric attacks because decision tree split optimization greedily minimizes overall impurity, allowing massive incoming novel classes to overwrite internal feature splits corresponding to minority historical classes.
    \item \textbf{Answer to RQ5 (Dynamic Allocation Efficacy):} Dynamically scaling replay quotas proportionally to incoming novelty volume restores class parity during tree induction, successfully capping maximum forgetting below $0.25$ across all evaluated high-volume attack vectors.
\end{itemize}

\section{Validation of Research Hypotheses}
The empirical results provide unequivocal validation for all four stated research hypotheses:
\begin{itemize}
    \item \textbf{Hypothesis H1 (Confirmed):} Baseline closed-world models yielded $F1 = 0.0000$ on all unobserved attack vectors.
    \item \textbf{Hypothesis H2 (Confirmed):} Isolation Forest provided zero-day screening alerts but exhibited insufficient recall ($18.86\%$) for standalone deployment.
    \item \textbf{Hypothesis H3 (Confirmed):} Memory overhead was reduced by $99.3\%$ without causing catastrophic legacy forgetting when maintaining a 100-sample/class floor.
    \item \textbf{Hypothesis H4 (Confirmed):} Volume-aware dynamic allocation structurally suppressed maximum forgetting spikes on volumetric attacks by $35.8\%$ relative to fixed static buffers.
\end{itemize}

\section{Limitations of the Present Work}
While demonstrating decisive improvements, several limitations characterize the present study:
\begin{enumerate}
    \item \textbf{Dataset Scope:} Evaluation was conducted primarily on the TON\_IoT benchmark. While highly representative of modern IoT telemetry, cross-dataset validation across other benchmarks (e.g., CIC-IoT-2023, Edge-IIoTset) would provide additional empirical generalization.
    \item \textbf{Uniform Exemplar Selection:} The dynamic allocation rule currently selects historical exemplars via stratified random sampling. While statistically sound, uniform sampling does not prioritize boundary-critical samples near decision margins.
    \item \textbf{Labeling Latency Assumption:} The pipeline assumes that out-of-distribution traffic flagged by Stage 1 is verified and assigned ground-truth labels by network telemetry or analysts prior to Stage 2 adaptation. In fully autonomous autonomous deployments, active learning or few-shot clustering would be required to minimize human labeling overhead.
\end{enumerate}

\section{Future Research Directions}
Several promising avenues for extending this research include:
\begin{enumerate}
    \item \textbf{Coreset and Boundary-Aware Exemplar Selection:} Integrating geometric coreset selection or uncertainty-based margin sampling to populate the dynamic replay buffer with hard boundary exemplars, further pushing catastrophic forgetting toward zero.
    \item \textbf{Federated Continual Learning across Edge Swarms:} Extending the dynamic replay mechanism to distributed federated architectures, enabling multiple IoT edge nodes to share localized threat exemplars without transmitting raw sensory telemetry over the network.
    \item \textbf{Hardware Acceleration on Edge Microcontrollers:} Deploying and profiling the continual learning engine on physical bare-metal microcontrollers (e.g., ARM Cortex-M, ESP32, RISC-V edge cores) to evaluate hardware energy dissipation and battery longevity under sustained zero-day attack campaigns.
\end{enumerate}
